\documentclass[10pt,a4paper]{article}
\usepackage[T1]{fontenc}
\usepackage[utf8]{inputenc}
\usepackage[dutch]{babel}
\usepackage{lmodern}
\usepackage{amsmath,amssymb,booktabs,graphicx,xcolor,geometry,hyperref,microtype}
\geometry{margin=21mm}
\hypersetup{colorlinks=true,linkcolor=blue,urlcolor=blue,citecolor=blue}
\newcommand{\vswirl}{\mathbf{v}_{\!\boldsymbol{\circlearrowleft}}}
\newcommand{\rhoF}{\rho_{\!f}}
\newcommand{\rhoE}{\rho_{\!E}}
\newcommand{\rhoM}{\rho_{\!m}}
\newcommand{\st}{\textsf}
\newcommand{\figroot}{../figures}
\title{\textbf{SST External Wave-Mode Benchmark v0.1.0}\\\large Bron-native fysica, reproduceerbare analytische controles en blind-gate ontwerp}
\author{SST-Workbench -- onafhankelijke literatuurbenchmark}
\date{19 oktober 2026\\\small Status: DEVELOPMENT / LITERATURE BASELINE / NOT CERTIFIED}
\begin{document}
\maketitle
\begin{abstract}
Twee door de gebruiker aangeleverde wetenschappelijke artikelen zijn geregistreerd als onafhankelijke referenties voor Kelvin-golfturbulentie en inertiële golven. De twee bron-PDF's worden met SHA-256 verzegeld. We hebben uitsluitend op de gepubliceerde parameters gebaseerde analytische tabellen gemaakt, plus negen geslaagde regressietests. Er zijn \emph{geen} experimentele punten gedigitaliseerd, geen ruwe PIV-data verwerkt en geen PKLSA-berekeningen uitgevoerd. De empirische validatiegates blijven \texttt{BLOCKED} en de SST-claim blijft \texttt{UNRESOLVED}.
\end{abstract}

\section{Bronnen, fysische regimes en waarneembare grootheden}
\begin{tabular}{p{.26\linewidth}p{.32\linewidth}p{.33\linewidth}}
\toprule
Bron & Bron-eigen observabele & Interpretatiebeperking\\ \midrule
Barckicke et al. (2026), PRL \cite{Barckicke2026PRL} & Lateral vortexkernverplaatsing $\eta(z,t)$; $S_\eta(k,\omega)$, spectrale hellingen, pentacoherentie. & Macro-waterwervel met luchtkern, continu gepompt. Kelvin-modi zijn niet gelijk aan SST-torsie.\\
Jackson--Thomas (2026), JFM \cite{JacksonThomas2026} & PIV-snelheidsvelden van vortexringomgeving, frequentie en golfrichting $\theta$. & Inertiële golven in bijna uniform roterende achtergrond. Uitsluitend een zijaanzicht van belletjes is geen PIV-veld.\\
\bottomrule
\end{tabular}

De PRL-bron vermeldt op PDF-p.~5 dat zijn bronmetingen \emph{niet publiek} zijn en op redelijk verzoek verkrijgbaar zijn. De JFM-bron biedt aanvullende video's, maar de ruwe PIV-cubes (orde TB) zitten niet in de meegeleverde PDFs. Het eerdere, aparte Nature Physics-onderzoek \cite{Barckicke2026Nature} verwijst naar een mogelijk bruikbare Figshare-dataset, DOI 10.6084/m9.figshare.30903347; deze is in v0.1.0 \emph{niet} ingelezen.

\section{Kelvin-modus: gepubliceerde parameters en een expliciet alternatieve kernbeschrijving}
Voor figuur~2 van \cite{Barckicke2026PRL} gelden de gerapporteerde waarden:
\begin{align}
a_0&=1.47\times10^{-3}\ \mathrm m,&\Gamma&=0.018\ \mathrm{m^2\,s^{-1}},&v_z&=-0.63\ \mathrm{m\,s^{-1}},\\
\Omega_0&=\frac{\Gamma}{2\pi a_0^2}=1325.739\ \mathrm{rad\,s^{-1}}.
\end{align}
Let op: figuren 3 en 4 in dezelfde paper gebruiken $a_0=1.3$~mm. Die twee geometrieën worden \emph{niet} samengevoegd. Voor de reproduceerbare analytische controle gebruiken we de \emph{hollow-core} benadering uit de aanvullende vergelijking~(7), en niet de impliciete Rankine-dispersie uit vergelijking~(1):
\begin{equation}
\omega_{\rm hc}(k,m)=\frac{\Gamma}{2\pi a_0^2}\left[1-\sqrt{1+x\frac{K_{m-1}(x)}{K_m(x)}}\right],\qquad x=a_0|k|,\quad m\ge 1.
\end{equation}
Hier is $K_m$ de gemodificeerde Besselfunctie van de tweede soort. Voor $m=1$ en $ka_0\ll1$ volgt de geverifieerde limiet
\begin{equation}
\omega_{\rm hc}\simeq -\frac{\Gamma k^2}{4\pi}\left[\ln\frac{2}{ka_0}-\gamma_{\rm E}\right].
\end{equation}
Dimensiecontrole: $[\Gamma k^2]=\mathrm{s^{-1}}$, evenals $[\Omega_0]$. Het teken volgt de Fourier-conventie van de bron; het is geen zelfstandig bewijs voor de voortplantingsrichting in een experiment.

\begin{center}\includegraphics[width=.93\linewidth]{\figroot/kelvin_reference.png}\end{center}
\noindent\textbf{Interpretatie:} de bovenstaande curve is \texttt{THEORY\_ONLY}; het is geen gedigitaliseerde experimentele dispersierelatie en geen PKLSA-voorspelling. De verticale advectie $v_z$ zit niet in deze versimpelde hollow-core uitdrukking. Dit is dus \emph{geen} volledige fit van figuur~2.

\section{Zes-golfturbulentie: falsifieerbare uitkomst, maar nog zonder brondata}
De bron rapporteert voor de golflengtespectra twee concurrerende voorspellingen:
\begin{align}
S_{\eta}^{\rm KS}(k)&\propto k^{-17/5},&S_{\eta}^{\rm KS}(\omega)&\propto\omega^{-11/5},\\
S_{\eta}^{\rm LN}(k)&\propto k^{-11/3},&S_{\eta}^{\rm LN}(\omega)&\propto\omega^{-7/3}.
\end{align}
De auteurs constateren expliciet dat de beide spectra \emph{te dicht bij elkaar} liggen om in hun metingen op basis van alleen de spectrale hellingen te onderscheiden. De zes-golfclaim wordt ondersteund met pentacoherentie op een 75 minuten durende meetreeks, niet alleen door power-law fits \cite{Barckicke2026PRL}. Zonder de complexe Fourier-amplitudes $\eta_k(t)$ kunnen we die onafhankelijke nulhypothesetoets niet reproduceren.
\begin{center}\includegraphics[width=.87\linewidth]{\figroot/cascade_exponents.png}\end{center}

\section{Inertiële golf: onafhankelijke standaardreferentie}
Voor bijna uniform roterende vloeistof geeft \cite{JacksonThomas2026} vergelijking~(1.2):
\begin{equation}
\omega_{\rm I}=\pm 2\,\Omega\cos\theta,\qquad \Omega=\frac{2\pi\,\mathrm{rpm}}{60},\qquad f_{\max}=\frac{2\Omega}{2\pi}.
\end{equation}
Dit is een \emph{intrinsieke} frequentie in het meedraaiende stelsel; in een stilstaande camera kunnen doppler- en azimutale mode-effecten optreden. Bij 12 rpm volgt $\Omega=1.25664\ \mathrm{rad\,s^{-1}}$ en $f_{\max}=0.400\ \mathrm{Hz}$. Bij $\theta=60^\circ$ volgt $f=0.200\ \mathrm{Hz}$. De bron vindt de nauwste overeenkomst met de dispersiecurve bij kleine Rossbygetallen (figuren 10 en 11), en rapporteert dezelfde relatie na onafhankelijke verwerking van axiale en radiale snelheidscomponenten.
\begin{center}\includegraphics[width=.93\linewidth]{\figroot/inertial_reference.png}\end{center}

\section{Verzegelde gates en onafhankelijk vergelijkingscontract}
\small
\begin{tabular}{lll}
\toprule
Gate & Status v0.1.0 & Betekenis\\ \midrule
G0 & PASS & Bronnen en SHA-256 vastgelegd\\
G1 & PASS & Brongrootheden, context en SI-consistentie\\
G2 & PASS & Hollow-core Eq.~(7), $ka_0\!\ll\!1$ regressie\\
G3 & PASS & Inertieel $\omega=2\Omega\cos\theta$ regressie\\
G4 & BLOCKED & Geen Kelvin-$\eta(z,t)$ of onafhankelijke meetpunten\\
G5 & BLOCKED & Geen ruwe spectrale onzekerheden voor hellingfit\\
G6 & BLOCKED & Geen 75-minuten amplitudes voor pentacoherentie\\
G7 & BLOCKED & Geen bron-native PIV-snelheidscubes\\
G8 & NOT RUN & Geen blind verzegelde PKLSA-predictor\\
G9 & UNRESOLVED & Fysische SST-hypothese niet beoordeeld\\
\bottomrule
\end{tabular}
\normalsize

De vooraf vastgelegde meetgates vereisen: onafhankelijke bronmetadata, correcte referentiestelseltransformatie, dezelfde observabelen, objectieve holdouts, geen fit van $\Gamma$ of $a_0$ aan de dispersiedata, en zo mogelijk een 95\%-betrouwbaarheidsinterval van de gerapporteerde onzekerheid. Een model slaagt niet doordat het zijn \emph{eigen} gegenereerde data reproduceert. Voor zes-golfinteracties wordt een fase-gerandomiseerde nulcontrole vereist met FDR $\le0.01$. Deze thresholds zijn \emph{preregistratieconcepten}; zonder ruwe meetonzekerheden nog niet toepasbaar.

\section{SST/PKLSA: minimale volgende stap}
\begin{enumerate}
\item Verkrijg $\eta(z,t)$ of de geregistreerde $S_\eta(k,\omega)$ van de Kelvin-golfmetingen en $u_x,u_r$ (plus co\"ordinaten, tijd en rotatiesnelheden) van de PIV-metingen; verifieer hashes en units.
\item Bevries de bron-native geometrie en randvoorwaarden. Een experimentele luchtkern op millimeterschaal mag niet stilzwijgend geïdentificeerd worden met de SST-canonieke $r_c=1.40897017\times10^{-15}$~m.
\item Voer een onafhankelijke PKLSA-vortexberekening uit die $\eta$, $\omega(k)$, $m$ en $v_z$ \emph{als output} levert. Geen Swirl-Clock-tijddilatiewet als gedwongen frequentie-input.
\item Vergelijk eerst met Euler/Rankine en inertiële golftheorie, dan pas met SST, inclusief modelverschil en onzekerheidsmarges.
\end{enumerate}

\paragraph{Begrensde conclusie.} De literatuur biedt reeds twee sterke, toetsbare referentieregimes. Het huidige product is een uitvoerbare \emph{literatuurbaseline}, geen SST-validatie en geen afgeronde blinde falsifier. De negen regressietests gaan uitsluitend over bronintegriteit en bekende mathematische limieten.

\begin{thebibliography}{9}
\bibitem{Barckicke2026PRL} J. Barckicke, C. Gissinger, and E. Falcon (2026), ``Experimental Evidence of Kelvin-Wave Turbulence along a Vortex Core'', \emph{Physical Review Letters} \textbf{137}, 134001. DOI: \href{https://doi.org/10.1103/t3bt-m431}{10.1103/t3bt-m431}; preprint \href{https://arxiv.org/abs/2607.07535}{arXiv:2607.07535v1}.
\bibitem{JacksonThomas2026} O. C. Jackson and P. J. Thomas (2026), ``An experimental study investigating the inertial-wave field accompanying vortex rings translating through a rotating fluid'', \emph{Journal of Fluid Mechanics} \textbf{1029}, A47. DOI: \href{https://doi.org/10.1017/jfm.2026.11199}{10.1017/jfm.2026.11199}.
\bibitem{Barckicke2026Nature} J. Barckicke, E. Falcon, and C. Gissinger (2026), ``Kelvin wave propagation along vortex cores'', \emph{Nature Physics} \textbf{22}, 409--414. DOI: \href{https://doi.org/10.1038/s41567-026-03175-w}{10.1038/s41567-026-03175-w}. Dataset: \href{https://doi.org/10.6084/m9.figshare.30903347}{10.6084/m9.figshare.30903347}.
\bibitem{Saffman1992} P. G. Saffman (1992), \emph{Vortex Dynamics}, Cambridge University Press. DOI: \href{https://doi.org/10.1017/CBO9780511624063}{10.1017/CBO9780511624063}.
\bibitem{Kelvin1880} W. Thomson (Lord Kelvin) (1880), ``Vibrations of a columnar vortex'', \emph{Philosophical Magazine} \textbf{10}, 155--168. DOI: \href{https://doi.org/10.1080/14786448008626912}{10.1080/14786448008626912}.
\end{thebibliography}
\end{document}
